\begin{table}[H]
\centering
\caption{Standardized Effect Sizes for Main Outcomes}
\label{tab:sde}
\begin{adjustbox}{max width=\textwidth}
\begin{tabular}{lccccccl}
\toprule
Outcome & $\hat{\beta}$ & SE & SD($X$) & SD($Y$) & SDE & SE(SDE) & Classification \\
\midrule
Inspections (90d post) & 0.0153 & 0.0293 & --- & 0.838 & 0.0183 & 0.0350 & Small positive \\
Violations (90d post) & -0.0180 & 0.0504 & --- & 1.421 & -0.0126 & 0.0355 & Small negative \\
Penalties (90d post) & -0.0275 & 0.1183 & --- & 3.301 & -0.0083 & 0.0358 & Small negative \\
\bottomrule
\end{tabular}
\end{adjustbox}
\par\vspace{0.3em}
{\footnotesize \emph{Notes:} This table reports standardized effect sizes (SDE) to facilitate cross-study comparison of treatment effect magnitudes. The treatment variable (news competition) is the standardized log of weekly FEMA disaster declarations, so SD($X$) = 1 by construction (marked ``---''). SDE $= \hat{\beta} / \text{SD}(Y)$, measuring the effect of a one-standard-deviation increase in news competition on the outcome in standard deviation units. \textbf{Research question:} Does news competition reduce post-fatality mine safety enforcement? \textbf{Treatment:} Continuous (standardized log weekly FEMA disaster declarations). \textbf{Data:} MSHA accidents, inspections, violations; FEMA disaster declarations; 1,069 fatality events, 2000--2024. \textbf{Method:} OLS reduced form with year-quarter and state fixed effects, heteroskedasticity-robust standard errors. \textbf{Sample:} All mine fatalities with matching FEMA weekly disaster data. Classification labels refer to the magnitude of the standardized point estimate, not to statistical significance. ``Null'' denotes a near-zero effect size ($|$SDE$| < 0.005$), not a failure to reject a null hypothesis.}
\end{table}

